% ============================================================
%  OR gate - transistor-level realization
%      OR = NOR followed by a NOT gate
%  Stage 1 (Q1,Q2 in parallel) = NOR  -> gives NOT(A+B)
%  Stage 2 (Q3)                = NOT   -> inverts back to (A+B)
%  Compile:  pdflatex circuit.tex
% ============================================================
\documentclass[border=12pt]{standalone}
\usepackage{circuitikz}
\usepackage{amsmath}
\usetikzlibrary{decorations.pathreplacing}

\begin{document}
\begin{circuitikz}[american, line width=0.7pt]

  % ---------- Supply rail (+5 V) ----------
  \draw (-2.5,5.5) -- (13,5.5);
  \node[anchor=south] at (5,5.55) {\textbf{$V_{CC} = +5\,\mathrm{V}$}};

  % ===================== STAGE 1 : NOR =====================
  \node[npn] (q1) at (3,2.5) {};
  \node[npn] (q2) at (6,2.5) {};
  \node[anchor=west] at ($(q1.E)+(0.4,0.15)$) {Q1};
  \node[anchor=west] at ($(q2.E)+(0.4,0.15)$) {Q2};

  % collectors tied, shared pull-up
  \draw (q1.C) -- (q2.C);
  \coordinate (top) at ($(q1.C)!0.5!(q2.C)$);
  \draw (top) to[R, l_={$R_{C1}=1\,\mathrm{k}\Omega$}] (top|-0,5.5);
  \fill (top|-0,5.5) circle (1.6pt);

  % emitters to ground
  \draw (q1.E) -- ++(0,-1.1) node[ground]{};
  \draw (q2.E) -- ++(0,-1.1) node[ground]{};

  % input A
  \draw (-2.5,2.5) node[ocirc]{} node[left]{Input $A$}
        to[R, l={$R_{B1}=10\,\mathrm{k}\Omega$}] (q1.B);
  % input B (underneath, vertical base resistor)
  \draw (-2.5,0.3) node[ocirc]{} node[left]{Input $B$}
        -- (4.6,0.3) to[R, l_={$R_{B2}=10\,\mathrm{k}\Omega$}] (4.6,2.5) -- (q2.B);

  % NOR node label
  \node[anchor=south] at (7.6,3.1) {$\overline{A+B}$};

  % ===================== STAGE 2 : NOT =====================
  \node[npn] (q3) at (10.5,2.5) {};
  \node[anchor=west] at ($(q3.E)+(0.4,0.15)$) {Q3};

  % interstage: NOR collector node -> base of Q3
  \draw (q2.C) -- (7.6,3.0) -- (7.6,2.5) to[R, l={$R_{B3}=10\,\mathrm{k}\Omega$}] (q3.B);

  % Q3 emitter to ground, collector pull-up
  \draw (q3.E) -- ++(0,-1.1) node[ground]{};
  \draw (q3.C) to[R, l_={$R_{C3}=1\,\mathrm{k}\Omega$}] (10.5,5.5);
  \fill (10.5,5.5) circle (1.6pt);

  % output
  \draw (q3.C) -- ++(2.2,0) node[ocirc]{} node[right]{Output $Y=A+B$};

  % ---------- stage brackets ----------
  \draw[decorate, decoration={brace, amplitude=5pt}, line width=0.5pt]
        (-2.6,-1.3) -- (7.0,-1.3) node[midway, below=6pt]{\small Stage 1: NOR};
  \draw[decorate, decoration={brace, amplitude=5pt}, line width=0.5pt]
        (8.2,-1.3) -- (12.9,-1.3) node[midway, below=6pt]{\small Stage 2: NOT};

\end{circuitikz}
\end{document}
